Lung adenocarcinoma (LUAD) progresses from atypical adenomatous hyperplasia (AAH) through
adenocarcinoma \emph{in situ} and minimally invasive adenocarcinoma to invasive disease.
Most transcriptomic descriptions of this sequence use bulk tissue, or tumour tissue alone.
We profiled paired single-nucleus RNA sequencing (413,697 nuclei, 23 patients) and Visium
spatial transcriptomics (56 sections, 25 patients) across all five stages, asking how the
transcriptional state of the tissue changes as invasion proceeds and whether it can be
reversed \emph{in silico} by known perturbations.

The tissue resolves into seven recurring niche-domain archetypes ($K^{*}=7$), six of them
ordinary lung structures and only one restricted to invasive disease. A confound then shaped
the analysis: sequencing depth and tissue density are the same measurement here, and domain
identity, marker-gene stability, the stage-dependence of the spatial programmes, the
cohort-level copy-number difference and one domain's prognostic weight all scale with it.
Several apparently positive results did not survive depth matching. The invasive domain's
ECM/interstitial identity did: the fibroblast ECM programme is higher in invasive disease in
20 of 22 paired cases ($P=0.0033$).

Querying a patient-paired disease axis against the LINCS library with three scoring schemes
returned 34 compounds in common, dominated by PI3K/mTOR inhibitors (split-half
$\rho=0.962$; cross-modality $\rho=0.735$). A positive control, recovery of the known
regulators of an official benchmark (TBXT 1st of 8,450), confirms that the pipeline recovers
true positives, so the negatives we report reflect the data rather than the method. Each is
reported with its diagnosis: no usable single-cell malignant/non-malignant separator, a
developmental trajectory running opposite to the source study, and no signal from gene-level
prediction.
